\section{From one percent agreement to list agreement} \label{sec:proof-of-lem-from-1-percent}
In this section we prove \pref{lem:from-1-percent-to-list-agreement}. 
\restatelemma{lem:from-1-percent-to-list-agreement}
Fix \(X\) to be a suitable simplicial complex and fix \(\mathcal{F} = \sett{f_r:r\to \Sigma}{r \in X(k)}\) to be an ensemble of functions on \(X(k)\). The distribution \(\mathcal{D}\) is \((\varepsilon_0,\eta)\)-sound. We also assume that \(\agr_{\mathcal{D}}(\mathcal{F})^{M}=\varepsilon^{M} \geq \varepsilon_0\) for some constant integer \(M > 0\) (which we do not explicitly calculate).

For a function \(g:X(0) \to \Sigma\) We denote by \(\supp_{\zeta}(g) = \sett{r \in X(k)}{f_r \overset{1-\zeta}{\approx} g}\) and call this set the \(\zeta\)-support of \(g\). We also denote by \(A_\zeta(g) = \sett{\set{r_1,r_2}}{r_1,r_2 \in \supp_\zeta (g) \ve f_{r_1}=f_{r_2}}\) the ``agreeing edge set of \(g\)'' and by \(\alpha_{\zeta}(g) = \Prob[\set{r_1,r_2} \sim D]{\set{r_1,r_2} \in A_{\zeta}(g)}\). Here \(D\) is the agreement distribution. When the ensemble of functions is not clear from context we denote \(\supp_{\zeta}^{\mathcal{F}}(g), A_{\zeta}^{\mathcal{F}}(g)\) etc.

For a face \(t \in X\), we denote by \(\mathcal{F}_t = \sett{f_r \in \mathcal{F}}{r \subseteq t}\). The set of functions whose agreeing set of edges is large is denoted by \(\overline{L}_{\tau,\zeta}(t) = \sett{g:t \to \Sigma}{\alpha_{\zeta}^{\mathcal{F}_t} (g) \geq \tau}\). In a lot of the following section we will consider partial functions \(g:t \to \Sigma\) (instead of all \(X(0)\)). In this case when we write \(\supp_\zeta(g), \alpha_\zeta(g)\) or \(A_\zeta (g)\) we will mean the support or agreement of \(g\) on its domain with respect to \(\mathcal{F}_t\).

\subsection{Overview - constructing the permutations in \pref{lem:from-1-percent-to-list-agreement}}
The idea is to use the the agreement theorem at hand to create lists of functions \(\overline{L}(s)=\set{L_s^1,L_s^2,\dots,L_s^\ell} \subseteq \overline{L}_{\tau,\zeta}(s)\) for every face \(s\). 
Then one needs to show that there is a matching between most lists \(\overline{L}(s), \overline{L}(t)\) such that \(s \subseteq t\). The matching we construct is so that \(\pi_{s,t}(i)=j\) if and only if \(L_t^i|_s \overset{1-\gamma}{\approx} L_t^j\) for some small \(\gamma>0\).

There are three degrees of freedom here, namely the choice of \(\tau,\zeta\) and \(\gamma\), and even after choosing them, it is not clear a priori which functions should appear in the list. We need the size of the lists to be \(\poly(\frac{1}{\varepsilon})\), so we cannot take all functions in \(\overline{L}_{\tau,\zeta}^0 (s)\). This gives rise to many problematic issues. The first two things we need to consider are:
\begin{enumerate}
    \item If \(f_r \overset{1-\gamma_1}{\approx} g|_r\) and \(g|_r \overset{1-\gamma_2}{\approx} g'|_r\) then \(f_r \overset{1-\gamma_1-\gamma_2}{\approx} g'|_r\). Hence, if two functions \(g,g':t \to \Sigma\) are close (in Hamming distance), then their support will have a lot of intersection. In particular, if \(L_t^i \in \overline{L}(t)\) we should exclude from the list a small Hamming ball around \(L_t^i\). Otherwise, if \(L_t^i|_s\) is \(\gamma\)-close to both \(L_s^j,L_s^{j'}\) we will have a difficulty of choosing whether \(\pi_{s,t}(i)=j\) or \(\pi_{s,t}(i)=j'\) in a way that will ensure that \(\pi_{s,u} = \pi_{s,t} \circ \pi_{t,u}\) for most \(s \subseteq t \subseteq u\).
    \item We want the lists to be exhaustive - meaning that every function \(g\) that has non-trivial agreement with the \(f_r\)'s has a close by function in the list. That is, if \(g:s \to \Sigma\) has that \(f_r \approx g|_r\) for \(\geq \tau=\poly(\varepsilon)\) fraction of the \(r \in X(k)\) inside \(s\), then either \(g \in \overline{L}(s)\) or \(g' \in \bar{L}(s)\) for some \(g' \approx g\). Otherwise, this could also lead to a difference in \(\overline{L}(s),\overline{L}(t)\)'s sizes, since for instance, it may be that \(g \in \overline{L}(t)\) but its restriction \(g|_s\) is not close to any element in the list of \(\overline{L}(s)\).
\end{enumerate}
Towards dealing with these issues we define \emph{density} and \emph{separation} of a set.
\begin{definition}
    Let \(\overline{L}_1,\overline{L}_2 \subseteq \set{g:s \to \Sigma}\). Let \(\gamma > 0\).
    \begin{enumerate}
        \item We say that \(\overline{L}_1\) is \(\gamma\)-separated if for every two \(g_1,g_2 \in \overline{L}_1\) it holds that \(\dist(g_1,g_2) > \gamma\).
        \item We say that \(\overline{L}_1\) is \(\gamma\)-dense in \(\overline{L}_2\) if for every \(h \in \overline{L}_2\) there exists \(g \in \overline{L}_1\) such that \(\dist(g,h) \leq \gamma\) (or as we denoted \(g \overset{1-\gamma}{\approx} h\)).
    \end{enumerate}
\end{definition}
Thus, for every \(s\) we will need to find a list \(\overline{L}(s)\) that is both separated (to overcome the first issue) and dense inside \(\overline{L}_{\tau,\zeta}(s)\) (to overcome the second issue). In fact, a method by \cite{DinurHKLT2018} allows us to find lists that are \(\gamma\)-dense but \(10\gamma\)-separated (for a smartly chosen \(\gamma\)). This is the content of \pref{claim:good-gamma-for-every-t-and-delta}. We will need this property in order to match between lists. If \(\overline{L}(s)\) is \(\gamma\)-dense in \(\overline{L}_{\tau,\zeta}(s)\) and \(\overline{L}(t)|_s = \sett{L_t^i|_s}{L_t^i \in \overline{L}(t)}\) is contained in \(\overline{L}_{\tau,\zeta}(s)\) then for every \(L_t^i \in \overline{L}(t)\) there exists some \(L_s^j \in \overline{L}(s)\) such that \(L_s^j \overset{1-\gamma}{\approx} L_t^i|_s\).\footnote{There is a subtlety here that even if \(\alpha_{\zeta}(L_t^i|_s) \geq \tau\), it could be that \(\alpha_{\zeta}(L_t^i) < \tau\). We will explain in the actual proof how to overcome this using sampling arguments, but for this overview let us just ignore this issue and assume that \(\alpha_{\zeta}(L_t^i|_s) = \alpha_{\zeta}(L_t^i)\).} %

%
The \(10\gamma\)-separation will promise that there is only a single such \(L_s^j\) for every \(L_t^i\), because if \(L_s^{j},L_s^{j'}\) are both \(\gamma\)-close to \(L_t^i\) then they are \(2\gamma\)-close to one another, which is a contradiction to \(10\gamma\)-separation. If we have such lists this implies that \(\pi_{s,t}(i)=j \Leftrightarrow L_t^i|_s \overset{1-\gamma}{\approx} L_s^j\) is a well defined function.

By making sure that \(\overline{L}(t)\) is also \(10\gamma\)-separated, we get that this function \(\pi_{s,t}\) is injective for most \(s \subseteq t\). If \(L_t^i,L_t^{i'}\) are \(10\gamma\)-separated, then for most \(s \subseteq t\), it will hold that \(L_t^i|_s, L_t^{i'}|_s\) will still be more than \(2\gamma\)-far. Hence, both \(L_t^i|_s, L_t^{i'}|_s\) cannot be \(\gamma\)-close to the same \(L_s^j\) by the triangle inequality. Finally, we show in the proof that the separation of all lists will ensure that their size stays \(\poly(\frac{1}{\varepsilon})\). For more details on the necessary condition for constructing these injections, see \pref{claim:injective-function} and its proof.

\medskip
Up until now we have explained in high level how to construct injective functions from \(\overline{L}(t)\) to \(\overline{L}(s)\) for \(t \supseteq s\). Surjectivity requires more care because \(\pi_{s,t}\) may still not be surjective even if \(\overline{L}(s),\overline{L}(t)\) are dense and separated. It could be that there is a ``new'' function \(L_s^j \in \overline{L}_{\tau,\zeta}(s)\) that is not close to \emph{any} function in the \(L_t^i|_s\). Another way of saying this is that \(\overline{L}(t)|_s\) may not be dense in \(\overline{L}_{\tau,\zeta}(s)\). %

To overcome this we will first show a structural property. We show that for some fixed \(t\), if \(\pi_{s,t}\) is not surjective for a non-negligible fraction of the \(s \subseteq t\), this implies that there is some \(\zeta' = O(\zeta)\), and a function \(g:t \to \Sigma\) such that \(\alpha_{\zeta'}(g) \geq \tau - \poly(\varepsilon)\), that is \(\zeta'\)-far from all functions in \(\overline{L}_{\tau,\zeta}(t)\). Showing this is where the previous agreement theorems come into play; if it is true that \(\overline{L}_{\tau,\zeta}(t)\) is not dense for many of the \(s \subseteq t\), this means that even after rerandomizing \(f_r\) for \(r \in \bigcup_{g \in \overline{L}_{\tau,\zeta}(t)} \supp_{\zeta}(g)\), the ensemble will pass the agreement test with non-negligible probability. By the agreement soundness guarantee, there is a new function \(g\) such that \(A_\zeta (g) \geq \tau^3\). We then show that by scaling \(\zeta\) by a constant factor, it also holds that \(A_{\zeta'} (g) \geq \tau-\poly(\varepsilon)\) (this happens because \(g\) needs to agree with the functions \(f_r\) for \(r \notin \bigcup_{g \in \overline{L}_{\tau,\zeta}(t)} \supp_{\zeta}(g)\), since these are the functions that were not rerandomized, although this needs to be meticulously argued).

On the other hand (keeping \(t\) fixed), we will show that there exists \(\zeta, \tau\) such that there are no new functions like this: We take any sequence \(\set{(\tau_i,\zeta_i)}_{i=1}^m\) where \(\tau_1=\varepsilon^2, \tau_{i+1}= \tau_i - \varepsilon^{50}\) and \(\zeta_{i+1} = 20\zeta_i\). Suppose that for every such pair there is a ``new'' \(g\) such that \(\alpha_{\zeta_{i+1}}(g) \geq \tau_{i+1} \geq \varepsilon^{10}\) that is far from all functions in \(\overline{L}_{\tau_i,\zeta_i}(t)\). This will imply that \(\prob{A_{\zeta_i}(g_i) \setminus \bigcup_{j=1}^{i-1} A_{\zeta_j}(g)} \geq \varepsilon^{10}\). This is because if \(g_i,g_{j}\) are far away from one another then most \(h_r\) that are close to \(g_i\) will be far from \(g_{j}\). Thus 
\[m\varepsilon^{10}\leq \sum_{i=1}^m \Prob[r_1,r_2]{A_{\zeta_i}(g_i) \setminus \bigcup_{j=1}^{i-1} A_{\zeta_j}(g_j)} = \Prob[r_1,r_2]{\bigcup_{i=1}^m A_{\zeta_i}(g_i)} \leq 1.\] 
We conclude that \(m \leq \varepsilon^{-10}\), i.e. that in every taking a long enough such sequence will result in some \(\set{(\tau_i,\zeta_i)}\) where this doesn't happen, which implies that \(\overline{L}_{\tau_i,\zeta_i}(t)\) is dense inside most \(\overline{L}_{\tau_i,\zeta_i}(s)\). This argument is depicted in \pref{claim:good-tau-and-delta-for-every-t}.

We use the pigeonhole principle to find a pair \((\tau_i,\zeta_i)\) and \(M_{(\tau_i,\zeta_i)} \subseteq X(d)\) of relative size \(\poly(\varepsilon)\), such that for every \(t \in M_{(\tau_i,\zeta_i)}\) and for all but a negligible fraction of \(s \subseteq t\) it holds that \(\overline{L}_{\tau_i,\zeta_i}(t)|_s\) is dense inside \(\overline{L}_{\tau_{i+1},\zeta_{i+1}}(s)\). Finally, using this set, we show that if \(t \in M_{(\tau_i,\zeta_i)}\) then most \(s \subseteq t\) also have the same property (i.e. that \(\overline{L}_{\tau_i,\zeta_i}(s)|_u\) is dense inside most \(\overline{L}_{\tau_{i+1},\zeta_{i+1}}(u)\) for most \(u \subseteq s\)). Using sampling, we conclude that this property propagates, i.e. that for a small enough dimension \(d' \ll d\) this density property holds for most \(t' \in X(d')\). This argument is in \pref{prop:good-parameters}.

Thus to conclude, we find our lists \(\overline{L}(s)\) by first finding the aforementioned \((\tau_i,\zeta_i)\). Then we find lists \(\overline{L}(s)\) that are \(\gamma\)-dense and \(10\gamma\)-separated inside \(\overline{L}_{\tau_{i+1},\zeta_{i+1}}(s)\) and define the permutation \(\pi_{s,t}(i)=j\) if and only if \(L_t^i|_s \overset{1-\gamma}{\approx} L_s^j\). These are well defined for most \(s \subseteq t\), and we use these for showing \pref{lem:from-1-percent-to-list-agreement}.

\subsection{Proof of \pref{lem:from-1-percent-to-list-agreement}}
This definition will be useful for the rest of this section.
\begin{definition}
    Let \(s \subseteq t\), \(\zeta > 0\) and \(L\) be a list of functions on \(t\). %
    We say that \(s\) \emph{samples \(L\) well} if %
    for every \(g_1,g_2 \in L\) it holds that \(\Abs{\dist(g_1|_s,g_2|_s) - \dist(g_1,g_2)} \leq \varepsilon^{100}\). We say that \(s\) \emph{samples \(L\)'s agreement} well if \(s\) samples \(L\) well and in addition, for every \(g \in L\) it holds that \(\Abs{\prob{A_{\zeta}(g)}- \prob{A_{\zeta}(g|_s)}} \leq \varepsilon^{100}\).
\end{definition}

We also define the requirements we need from the lists \(\overline{L}(s)\).
\begin{definition} \label{def:mi}
    Let \(d_1 \leq d_2 \leq d_3 \leq d\) be dimensions. Let \(\tau,\zeta,\gamma > 0\) be constants and let \(\ell\) be integers. Let \(M_i = M(\tau,\zeta,\gamma,\ell) \subseteq X(d_i)\) be all the \(t \in X(d_i)\) such that there exists some \(\overline{L}(t) \subseteq \overline{L}_{\tau,\zeta}(t)\) such that:
    \begin{enumerate}
        \item \(\overline{L}(t)\) is \(9\gamma\)-separated.
        \item \(\overline{L}(t)\) is \(\gamma\)-dense in \(\overline{L}_{\tau - \varepsilon^{100},\zeta}(t)\).
        \item \(|\overline{L}(t)|=\ell\).
        \item For \(i=2,3\), and \(i'< i\). Then the following set \(P_t(i')\) has probability \(\Prob[s \subseteq t, s \in X(i')]{P_t(i')} \geq 1-\exp(-\Omega(\poly(\varepsilon)\frac{d_1}{k}))\). \(P_t(i')\) is all the \(s \subseteq t, s \in X(i')\) such that \(s\) samples \(\overline{L}(t)\)'s agreement well.
    \end{enumerate}    
\end{definition}

The proof of \pref{lem:from-1-percent-to-list-agreement} relies on the following proposition.
\begin{proposition} \label{prop:good-parameters}
    There exists parameters \(\tau = \poly(\varepsilon),\zeta,\gamma \leq \eta\exp(\poly(1/\varepsilon))\) and \(\ell=\poly(1/\varepsilon)\) such that the following holds. Let \(d_1\leq d_2 \leq d_3 \leq d\) be such that \(\frac{d_3}{d} = \exp(-\poly(\varepsilon)\frac{d_1}{k})\). Then \(\prob{M_i} \geq 1- \exp(-\poly(\varepsilon)\frac{d_1}{k}) \).
\end{proposition}

We will also use this claim,
\begin{claim} \label{claim:injective-function}
    Let \(L_1 = \set{g_1,g_2,\dots,g_{\ell_1}}, L_2=\set{h_1,h_2,\dots,h_{\ell_2}}\) be such that \(L_1\) is \(\gamma\)-dense in \(L_2\) and such both \(L_1\) and \(L_2\) are \(2\gamma\)-separated. Then there exists an injective function \(\pi:[\ell_1]\to [\ell_2]\) such that \(\pi(j)=k\) if and only if \(g_j \overset{1-\gamma}{\approx} h_k\).
\end{claim}

\begin{proof}[Proof of \pref{lem:from-1-percent-to-list-agreement}]
Recall the definition of \(M_i\) is \pref{def:mi}. Given \pref{prop:good-parameters}, for every \(s \in M_1, t \in M_2\) and \(u \in M_3\) we take such a list \(\overline{L}(x) = \set{L_x^1,L_x^2,\dots,L_x^\ell}\), where \(x\) is \(s,t,u\) respectively (for \(x \notin M_i\) we take arbitrary lists). By definition, for every \(x \in M_i\) and every function in \(\overline{L}_x\) satisfies item \((1)\) in \pref{lem:from-1-percent-to-list-agreement} by definition. Let us verify the second and third items. We show the following for pairs \(s \subseteq t\), \(s \subseteq u\) and \(t \subseteq u\), we denote such a pair as \(x \subseteq y\) for short. For such a pair where \(y \in M_i, x \in P_y(i') \cap M_{i'}\) (where \(i,i'\) are the corresponding sets) we define \(\pi_{x,y}:[\ell]\to [\ell]\) by
\begin{equation} \label{eq:def-of-permutation}
    \pi_{x,y}(j)=k \Leftrightarrow L_y^j|_x \overset{1-\gamma}{\approx} L_x^k.
\end{equation} 
If either \(y \notin M_i\) or \(x \notin P_y(i')\cap M_{i'}\) we take \(\pi_{s,t}=Id\) as an arbitrary choice. We show that the assumptions in \pref{claim:injective-function} hold, which implies that this function is well defined. If this permutation is well defined for \(1-\exp(-\poly(\varepsilon)\frac{d_1}{k})\) of the pairs \(x \subseteq y\) this implies item \((2\) in \pref{lem:from-1-percent-to-list-agreement}.

First we note that if \(x\) samples \(\overline{L}(y)\)'s agreement well then \(\overline{L}(y)|_x \subseteq \overline{L}_{\tau-\varepsilon^{100},\zeta}(x)\) and \(\overline{L}(y)|_u\) is \(8\gamma\)-separated. Moreover if \(\overline{L}(x)\) is \(\gamma\)-dense in \(\overline{L}_{\tau-\varepsilon^{100},\zeta}(x)\), and \(\overline{L}(y)|_x \subseteq \overline{L}_{\tau-\varepsilon^{100},\zeta}(x)\) then in particular \(\overline{L}(x)\) is \(\gamma\)-dense in \(\overline{L}(y)|_x\). Finally, by definition \(\overline{L}(y)\) is also \(9\gamma\)-separated. Hence \pref{claim:injective-function} implies that this is a well defined and injective function. As both \(\overline{L}(y)|_x\) and \(\overline{L}(y)\) have the same size \(\ell\), this is indeed a permutation.

Finally, we will prove the third item in \pref{lem:from-1-percent-to-list-agreement},
\begin{equation*}
    \Prob[s \subseteq t \subseteq u, s \in X(d_1),t \in X(d_2), u \in X(d_3)]{\pi_{s,t}\circ \pi_{t,u} = \pi_{s,u}} \geq 1-\exp(-\Omega(\poly(\varepsilon) \frac{d_1}{k}).
\end{equation*}
Note that the following events are all occur with probability \(1-\exp(-\Omega(\poly(\varepsilon) \frac{d_1}{k})\):
\begin{enumerate}
    \item \(u \in M_1, s \in M_2, t \in M_3\).
    \item \(t \in P_u(i_2), s \in P_t(i_1) \cap P_u(i_1)\).
    \item For every \(j\) and \(k\), 
    \begin{equation} \label{eq:dist-preserved-in-list-restriction}
        \Abs{\dist(L_t^j|_s,L_u^k|_s) - \dist(L_t^j,L_u^k|_t)} \leq \gamma.
    \end{equation}
\end{enumerate}
The first two items are by \pref{prop:good-parameters}. The last item is via the following claim, proven later on (we will need this claim also for later):
\begin{claim} \label{claim:most-sets-sample-well}
    Fix \(t\), \(L\) and \(\zeta\). Then the fraction of \(s \subseteq t, s\in X(d_1)\) that \emph{don't} sample \(L\)'s agreement well is at most \(|L|^2 \exp(-\poly(\varepsilon) \frac{d_1}{k})\).
\end{claim}
Indeed, if \(t\) samples \(\overline{L}(u)\) well and \(s\) samples \(\overline{L}(u)|_t \cup \overline{L}(t)\) well, then distances between members of \(\overline{L}(u)|_t \cup \overline{L}(t)\) are preserved and the third item holds (and both these events occur with probability  \(1-\exp(-\poly(\varepsilon) \frac{d_1}{k})\)).

Thus to prove the theorem we assume \(s \subseteq t \subseteq u\) are such that all these events occur and show that \(\pi_{s,t}\circ \pi_{t,u} = \pi_{s,u}\). Let \(k = \pi_{s,t}\circ \pi_{t,u}(j)\) and \(k'=\pi_{s,u}(j)\). Then it holds that \(L_u^j|_s \overset{1-\gamma}{\approx} L_s^{k'}\), and that \(L_u^j|_t \overset{1-\gamma}{\approx} L_t^{\pi_{t,u}(j)}\) and \(L_t^{\pi_{s,t}(\pi_{t,u}(j))}|_s \overset{1-\gamma}{\approx} L_s^{k}\). By \eqref{eq:dist-preserved-in-list-restriction}, distances are approximately preserved, so \(L_u^j|_s \overset{1-2\gamma}{\approx} L_t^{\pi_{t,u}(j)}|_s\). Using the triangle inequality (where \(L^{\pi_{t,u}(j)}_t|_s\) is the point in the middle) it holds that \(L_u^j|_s \overset{1-3\gamma}{\approx} L_s^{k}\), which implies that \(L_s^k \overset{1-4\gamma}{\approx} L_s^{k'}\). By \(10\gamma\)-separation of \(\overline{L}(u)\), \(k=k'\).
\end{proof}

The proof of \pref{claim:injective-function} is direct.
\begin{proof}[Proof of \pref{claim:injective-function}]
     By density of \(L_2\) in \(L_1\), it holds that for every \(g_j\) there exists some \(h_k\) such that \(g_j \overset{1-\gamma}{\approx} h_k\). Moreover, there is only one such \(k\): assume that for some \(j\) there are \(k,k'\) such that \(g_j \overset{1-\gamma}{\approx} h_k\) and \(g_j \overset{1-\gamma}{\approx} h_{k'}\). Then by the triangle inequality \(h_k \overset{1-2\gamma}{\approx} h_{k'}\), which by \(2\gamma\)-separation implies that \(k=k'\). Thus it is a well defined function.

The same argument shows that this is an injection. Indeed, Let \(k=\pi(j)=\pi(j')\). That is, \(h_j \overset{1-\gamma}{\approx} h_k \overset{1-\gamma}{\approx} h_{j'}\). Then \(h_j \overset{1-2\gamma}{\approx} h_{j'}\). by \(2\gamma\)-separation of \(\overline{L}(s)|_u\), this implies that \(j=j'\), i.e. the function is an injection.
\end{proof}

So is the proof of \pref{claim:most-sets-sample-well}.
\begin{proof}[Proof of \pref{claim:most-sets-sample-well}]
    Let us start with sampling well. There are \(\leq \binom{|L|}{2}\) pairs \(g,g' \in L\) so by a union bound it is enough to argue that for at most \(\exp(-\poly(\varepsilon)d_1)\) of the \(s \subseteq t\) it holds that \(\Abs{\dist(g,g')-\dist(g|_s,g'|_s)} > \varepsilon^{100}\). Let \(A \subseteq t\) be the set of \(v \in t\) such that \(g(v)\ne g'(v)\). By \pref{thm:sampling-in-HDXs} the fraction of \(s \subseteq t\) such that \(\Abs{\prob{A} - \Prob[v \subseteq s]{A}} > \varepsilon^{100}\) is \(\exp(-\poly(\varepsilon)d_1)\). On the other hand, \(\prob{A} = \dist(g,g')\) and \(\Prob[v \subseteq s]{A} = \dist(g|_s,g'|_s)\) giving us the that at most \(\binom{|L|}{2}\exp(-\poly(\varepsilon)d_1)\) don't sample \(L\) well.

    Now let us make sure that most \(s \subseteq t\) also sample \(L\)'s agreement well. Indeed, let \(g \in L\). By \pref{claim:edge-sampling} there is at most \(\exp(-\poly(\varepsilon)\frac{d_1}{k})\) of the \(s \subseteq t\) such that 
    \[\Abs{A_{\zeta}(g) - A_{\zeta}(g|_s)} > \varepsilon^{100}.\]
    Another union bound gives us us the claim.
\end{proof}

The proof of \pref{prop:good-parameters} will follow from these assertions. We prove them after the proof of the proposition.
\begin{claim} \torestate{\label{claim:good-tau-and-delta-for-every-t}
    Let \(d_1 \leq d\) and fix \(t \in X(d)\) such that \(\agr_{\mathcal{D}}(\mathcal{F}_t) \geq \varepsilon - \varepsilon^{100}\). There exists \(m\leq \frac{1}{\varepsilon^{10}}\) such that the following holds for \(\tau_m = \varepsilon^2(1-\varepsilon^{50})^m\) and \(\zeta_m = 3\eta 20^m\). There exists \(L \subseteq \overline{L}_{\tau_m,\zeta_m}(t)\) such that for \(1-\exp(-\poly(\varepsilon)\frac{d_1}{k})\) of the \(u \subseteq t, u \in X(d_1)\), \(L|_u\) is \(20\zeta_m\)-dense in \(\overline{L}_{\tau_{m+1},\zeta_m}(u)\).}
\end{claim}

\begin{claim} \label{claim:good-gamma-for-every-t-and-delta}
    Let \(\zeta,\tau > 0\) and let \(t \in X(d)\). Then there exists some \(\gamma = \frac{24}{23} 24^i \zeta\) for \(i=1,2,\dots,\varepsilon^{-10}\) and some \(L \subseteq \overline{L}_{\tau,\zeta}(t)\) such that \(L\) is \(\gamma\)-dense, \(\abs{L} \leq \frac{1}{\tau^3}\) and \(23\gamma\)-separated.
\end{claim}

\begin{proof}[Proof of \pref{prop:good-parameters}]
    By \pref{claim:edge-sampling} there are \(1-\exp(-\poly(\varepsilon)\frac{d_1}{k})\) faces \(t \in X(d)\) such that \(\agr_{\mathcal{D}}(\mathcal{F}_t) \geq \varepsilon - \varepsilon^{100}\) (we use \pref{claim:edge-sampling} on the set of \(\set{r_i}\) such that \(\set{f_{r_i}}\) pass the agreement test).
    By \pref{claim:good-tau-and-delta-for-every-t}, for every such \(t \in X(d)\) there exists \[\tau = \varepsilon^2(1-\varepsilon^{50})^i, \zeta = 3\eta 20^i\] (for \(i \in \set{1,2,\dots,\frac{1}{\varepsilon^{10}}}\)), and a list \(L'(t) \subseteq \overline{L}_{\tau,\zeta}(t)\) as in \pref{claim:good-tau-and-delta-for-every-t}. 
    
    Moreover, by \pref{claim:good-gamma-for-every-t-and-delta} there exists some 
    \[\gamma' \in [\frac{24}{23} 24\zeta,\frac{24}{23} 24^{\frac{1}{\varepsilon^{10}}}\zeta]\]
    and \(\overline{L}(t) \subseteq \overline{L}_{\tau,\zeta}(t)\) such that \(\overline{L}(t)\) has length \(\leq \poly(1/\varepsilon)\), \(\overline{L}(t)\) is \(\gamma'\)-dense in \(\overline{L}_{\tau,\zeta}(t)\) and \(23\gamma'\)-separated.

    Thus by the pigeonhole principle, there exists \(\tau,\zeta,\gamma',\ell\) and a set \(N_d = N_d(\tau,\zeta,\gamma',\ell) \subseteq X(d)\) of probability \(\prob{N_d} \geq \poly(\varepsilon)\), such that for every \(t \in X(d)\) there exists \(L'(t) \subseteq \overline{L}_{\tau,\zeta}(t)\) as in \pref{claim:good-gamma-for-every-t-and-delta} and \(\overline{L}(t) \subseteq \overline{L}_{\tau,\zeta}(t)\) of size \(\Abs{\overline{L}(t)}=\ell\), \(\overline{L}(t)\) is \(\gamma\)-dense in \(\overline{L}_{\tau,\zeta}(t)\) and \(23\gamma\)-separated.

    Let us define \(M_i \subseteq X(d_i)\) be the set of \(s \in X(d)\) such that there exists some \(t \in N_d, t\supset s\) such that:
    \begin{enumerate}
        \item \(s \subseteq t\) samples \(L'(t) \cup \overline{L}(t)\) well.
        \item \(L'(t)\) is \(20\zeta\)-dense in \(\overline{L}_{\tau-\varepsilon^{100},\zeta}(s)\).
        \item If \(i=2,3\) and \(i' < i\), then \(\Prob[u \subseteq s]{M_{i'}} \leq \exp(-\poly(\varepsilon)\frac{d_1}{k}))\) (This definition is recursive but \(M_i\) is always well defined).
    \end{enumerate}
    We show that for every \(s \in M_i\) the proposition holds for \(\tau,\zeta,\ell\) and \(\gamma=2\gamma'\). We take \(\overline{L}(s)=\overline{L}(t)|_s\) for some arbitrary \(t \in N_d\) that contains \(s\) and such that the items above holds for this \(t\). By the fact that \(\overline{L}(t)\) is sampled well we have that \(\overline{L}(s) \subseteq \overline{L}_{\tau-\varepsilon^{100},\zeta}(t)\). Moreover if \(\overline{L}(t)\) is \(23\gamma'=11.5\gamma\)-separated and in particular it holds that \(11\gamma'\)-separated (and in particular, all restrictions indeed result in distinct functions, i.e. \(\Abs{\overline{L}(s)}=\Abs{\overline{L}(t)|_s}=\ell\)). 
    
    Next we show density in \(\overline{L}_{\tau-\varepsilon^{100},\zeta}(s)\). Note that \(L'(t)|_s\) is \(20\zeta\)-dense in \(\overline{L}_{\tau-\varepsilon,\zeta}(s)\), and \(\overline{L}(t)\) is \(\gamma'\)-dense in \(L'(t)\) (not to say that this is a subset of \(L'(t)\), just that for every \(g' \in L'(t)\) there exists some \(g \in \overline{L}(t)\) such that \(g' \overset{1-\gamma'}{\approx} g\)). Thus by the fact that \(s\) samples \(L'(t) \cup \overline{L}(t)\) well, this implies that \(\overline{L}(t)\) is \(\gamma'+\varepsilon^{100}\)-dense in \(L'(t)\). Hence, it holds that \(\overline{L}(t)\) is \(\gamma'+\varepsilon^{100}+20\zeta\leq \gamma\)-dense in \(\overline{L}_{\tau-\varepsilon^{100},\zeta}(s)\).

    Finally, we show that when \(i=2,3\) and \(i' \leq i\), it holds that
    \(P_s(i')\) from the definition of the proposition has size \(1-\exp(-\Omega(\poly(\varepsilon)\frac{d_1}{k}))\). Let us go one by one:
    \begin{enumerate}
        \item For every \(t\), the fraction of \(u \subseteq s\) that don't sample \(\overline{L}(s)\) well is \(\exp(-\Omega(\poly(\varepsilon)\frac{d_1}{k}))\) by \pref{claim:most-sets-sample-well}. %
        \item The probability that \(u \subseteq s\) is in \(M_{i'}\) is  \(\exp(-\Omega(\poly(\varepsilon)\frac{d_1}{k}))\). Let us show that in this case \(\overline{L}(t)|_s\) is \(2\gamma\)-dense. \(u \in M_{i'}\), therefore there is some \(\overline{L}(u)\) of size \(\ell\) that is \(\gamma\)-dense in \(\overline{L}_{\tau-\varepsilon^{100},\zeta}(u)\), and \(11\gamma\)-separated (by what we already proved above). Thus if we show that for every \(g \in \overline{L}(u)\) there is some \(g' \in \overline{L}(t)|_s\) such that \(g \overset{1-\gamma}{\approx} g'\) it will follow that \(\overline{L}(t)|_s\) is \(2\gamma\)-dense. By \pref{claim:injective-function}, the fact that \(\overline{L}(u)\) is dense in \(\overline{L}(t)|_s\) and the fact that both are \(2\gamma\)-separated, there is an injective function \(\phi:\overline{L}(t)|_s \to \overline{L}(u)\) where \(\phi(g') = g\) if and only if \(g \overset{1-\gamma}{\approx} g'\). As both lists have the same size, this shows that this is a surjection, i.e. that for every \(g \in \overline{L}(u)\) there is some \(g' \in \overline{L}(t)|_s\).
    \end{enumerate}
    
    To conclude, let us show that \(\prob{M_i} \geq 1- \exp(-\Omega(\poly(\varepsilon)\frac{d_1}{k})\). Let us begin with \(M_1\). Let \(B_1\) be the fraction of \(s\)'s such that \(\Prob[t \supseteq s]{N_d} \leq \poly(\varepsilon)\). Its fraction is \(\prob{B_1} = O(\frac{d_1}{d\poly(\varepsilon)}) = \exp(-\poly(\varepsilon)\frac{d_1}{k})\) by \pref{cor:good-sampler-from-expansion}. Let \(B_2 \subseteq X(d_i)\) be the event that more than \(\frac{1}{3}\) of the \(t \supseteq s, t\in N_d\) have:
    \begin{enumerate}
        \item Either \(s\) doesn't sample \(\overline{L}(t) \cup L'(t)\) well.
        \item \(L'(t)\) is \(20\zeta\)-dense in \(\overline{L}_{\tau-\varepsilon^{100},\zeta}(s)\).
    \end{enumerate}
    On the one hand, by \pref{claim:most-sets-sample-well} and \pref{claim:good-tau-and-delta-for-every-t} there are at most \(\exp(-\poly(\varepsilon)\frac{d_1}{k})\) such pairs. On the other hand, every \(s \in B_2 \setminus B_1\) contributes \(\poly(\varepsilon)\) such pairs. By Markov's inequality, \(\prob{B_2 \setminus B_1} = \exp(-\poly(\varepsilon)\frac{d_1}{k})\). Hence, if \(s \notin B_1 \cup B_2 = B_1 \cup (B_2 \setminus B_1)\), then \(s \in M_1\) so \(\prob{M_1} =  1 - \exp(-\poly(\varepsilon)\frac{d_1}{k})\).

    Continuing with \(M_2\). The same arguments for \(M_1\) show that \(1-\exp(-\poly(\varepsilon)\frac{d_1}{k})\) of the \(s \in X(d_2)\) sample \(\overline{L}(t)\cup L'(t)\) well and that \(L'(t)\) is \(20\zeta\)-dense in \(\overline{L}_{\tau-\varepsilon^{100},\zeta}(s)\). So we concentrate on the last property. The probability that \(s \in X(d_2)\) contains more than \(\\prob{M_1}^{1/2}\)-fraction of \(u \notin M_1\) is at most \(\left ( \prob{M_1} \right )^{1/2}\) (by Markov's inequality), and this is also \(\exp(-\poly(\varepsilon)\frac{d_1}{k})\), albeit with a slightly worse \(\poly(\varepsilon)\) than before, since since \(\prob{M_1}=\exp(-\poly(\varepsilon)\frac{d_1}{k})\). Hence \(\prob{M_2} =1-\exp(-\poly(\varepsilon)\frac{d_1}{k})\). The proof for \(M_3\) is similar.
\end{proof}

\begin{proof}[Proof of \pref{claim:good-gamma-for-every-t-and-delta}]
    Consider the following method, resembling \cite{DinurHKLT2018}. Let \(L_1 \subseteq \overline{L}_{\tau,\zeta}(t)\) be a maximal \(24\zeta\)-separated list. Consider the following chain \(L_1 \supseteq L_2 \supseteq L_3 \dots\) where every \(L_i\) is a maximal \(24^i \zeta\)-separated set inside \(L_{i-1}\). Obviously, the sizes of these sets are monotonically decreasing (but they are always non-empty), therefore there must be some \(j\) such that \(L_j = L_{j+1}\). Take as \(L=L_i\) for \(i\) being the first such \(j\).

    Recall that \(\gamma_i = \frac{24}{23} 24^i\). By definition, \(L_{i+1}\) is \(23\gamma_i = 24^{i+1}\zeta\)-separated. Let us see that is it also \(26\cdot 24^i \zeta\)-dense in \(\overline{L}_{\tau,\zeta}(t)\). First note that every \(L_j\) is \(24^i \zeta\)-dense in \(L_{j-1}\) (if it weren't, we could have added to \(L_j\) another \(g \in L_{j-1}\setminus L_j\) that is far from the current \(L_j\), contradicting maximality). Thus for every \(g \in \overline{L}_{\tau,\zeta}(t)\) there is a sequence \(g=g_0,g_1,g_2,\dots, g_i\) where \(g_j \in L_j\) for \(j>0\), such that \(\dist(g_{j-1},g_j)\leq 24^j \zeta\). Thus concluding that
    \[\dist(g,g_i) \leq \sum_{j=1}^i 24^j  \zeta = \frac{24}{23}(24^i-1)\zeta \leq \frac{24}{23} \cdot 24^i = \gamma_i.\]

    Finally, we must show that \(i \leq \frac{1}{\tau^3}\). If we show that \(\Abs{L_1} \leq \frac{1}{\tau^3}\) this will follow, since \(\Abs{L_i}\geq 1\) and for \(j \leq i\), \(\Abs{L_j}\leq \Abs{L_{j-1}}-1\). We distinguish this part of the proof to a separate claim since we will need it in also later on.
    \begin{claim} \label{claim:lists-are-short}
        let \(\tau \gg \exp(-\Omega(\zeta^2 k))\) and let \(\zeta\) and \(\gamma \geq 3\zeta\). Let \(L \subseteq \overline{L}_{\tau,\zeta}(t)\) be a \(\gamma\)-separated set. Then \(\Abs{L}\leq\frac{1}{\tau^3}\).
    \end{claim}
\end{proof}

\begin{proof}[Proof of \pref{claim:lists-are-short}]
    Assume towards contradiction that \(|L_1|>\frac{1}{\tau^3}\) and take some subset \(L'\subseteq L_1\) of size \(\frac{1}{\tau^3}+1\). As it is at least \(3\zeta\)-separated, for every distinct \(g,g' \in L'\), the fraction of \(r \in X(k)\) such that \(\dist(g|_r,g'|_r) \leq 2\zeta\) (i.e. \(g|_r \overset{1-2\zeta}{\approx} g'|_r\)) is at most \(\exp(-\Omega(\zeta^2 k))\leq \tau^{10}\). This shows that there are almost no intersection between the supports of \(g\) and \(g'\), since \(\supp_\zeta (g) \cap \supp_{\zeta}(g') \subseteq \sett{r \in X(k)}{\dist(g|_r,g'|_r) \leq 2\zeta}\). That is, \(\prob{\supp_\zeta (g) \cap \supp_{\zeta}(g')} \leq \exp(-\Omega(\zeta^2 k))\leq \tau^{10}\).
    
    Denote by 
    \[D(g) = A_{\zeta}(g) \setminus \bigcup_{g' \in L', g' \ne g}A_{\zeta}(g').\]
    Then in particular 
    \begin{align}
        \prob{D(g)} &\geq \prob{A_{\zeta}(g)} - \sum_{g' \in L', g' \ne g}\prob{A_\zeta (g) \cap A_\zeta (g')} \\
        &\geq \prob{A_{\eta}(g)} - \sum_{g' \in L', g' \ne g} 2 \prob{\supp_{\zeta}(g) \cap \supp_{\zeta}(g')} \\
        &\geq \tau - 2\tau^6.
    \end{align}
    On the other hand, the sets \(\set{D(g)}_{g \in L'}\) are mutually disjoint, namely
    \(1 \geq \sum_{g \in L'} \prob{D(g)} \geq \frac{\tau-2\tau^6}{\tau^3+1}>1\) which is a contradiction.
\end{proof}

\subsection{Proof of \pref{claim:good-tau-and-delta-for-every-t}}
We start with a weaker claim than \pref{claim:good-tau-and-delta-for-every-t}, saying that \(\overline{L}_{\tau^3,\zeta}(t)\) restricted to \(s \subseteq t\) is dense in \(\overline{L}_{\tau,\zeta}(s)\) for most \(s\). This claim is where we use the agreement soundness of the original distribution.

\begin{claim} \label{claim:density-of-smaller-tau}
    Fix \(t\),\(\tau\), \(\zeta\) and \(\gamma \geq 3\zeta\). Let \(L \subseteq \overline{L}_{\tau^3,\zeta}(t)\) be a \(\gamma\)-dense and \(\gamma\)-separated list. Then
    \begin{equation*}
        \Prob[s \subseteq t,s \in X(d_1)]{L|_s \text{ is \((\gamma+4\zeta)\)-dense in } \overline{L}_{\tau,\zeta}(s)} \geq 1-\exp(-\Omega(\frac{d_1}{k})).
    \end{equation*}
\end{claim}

\begin{proof}[Proof of \pref{claim:density-of-smaller-tau}]
    The claim will follow from these assertions.
    \begin{enumerate}
        \item If \(L\) is \(\gamma\)-dense and has size \(\poly(1/\varepsilon)\) in \(\overline{L}_{\tau^3,\zeta}(t)\) then 
        \[\Prob[r_1,r_2 \subseteq t, r_1,r_2 \sim D]{f_{r_1}=f_{r_2} \ve r_1,r_2 \notin \bigcup_{g \in L}\supp_{\gamma+2\zeta}(g)} \leq 2\tau^{1.5}.\]
        \item This quantity is sampled well. That is, let \(P = \sett{(r_1,r_2)}{f_{r_1}=f_{r_2} \ve r_1,r_2 \notin \bigcup_{g \in L}\supp_{\gamma+2\zeta}(g)}\). Then for \(1-\exp(-\Omega(\poly(\varepsilon)\frac{d_1}{k}))\) of the \(s \subseteq t\) it holds that
        \[\Abs{\Prob[r_1,r_2 \subseteq s]{P} - \prob{P}} \leq \varepsilon^{100}.\]
    \end{enumerate}
    Let us explain why. If there exists some \(h \in \overline{L}_{\tau,\zeta}(s)\) such that \(\dist(h,g) \geq \gamma+4\zeta\) for all \(g \in L|_s\), then \(A_{\zeta}(h) \geq \tau\) and moreover it holds that \(\prob{\supp_\zeta(h) \cap \supp_{\gamma+2\zeta}(g)} = \exp(-\poly(\varepsilon)k) \ll \poly(\varepsilon)\) (since whenever \(f_r\overset{1-\zeta}{\approx}h|_r\) and \(f_r\overset{1-\gamma-2\zeta}{\approx}g|_r\) then it holds that \(\dist(g|_r,h|_r)\leq \gamma+3\zeta\) which happens with \(\exp(-\poly(\varepsilon)k)\) probability at most). By \pref{claim:lists-are-short}, \(|L|=\poly(1/\varepsilon)\) so all these intersections are negligible. In particular, more than (say) half of the edges \((r_1,r_2) \in A_{\zeta}(h)\) have that \(r_1,r_2 \notin \bigcup_{g \in L}\supp_{\gamma+3\zeta+\varepsilon^{100}}(g)\). Thus \(\Prob[r_1,r_2 \subseteq s, r_1,r_2 \sim D]{P} \geq \frac{1}{2} A_\zeta(h) \geq \frac{1}{2}\tau\). On the other hand, by the first item \(\prob{P} \leq 2\tau^{1.5}\). By the second item \(\Prob[r_1,r_2 \subseteq s]{P} \leq \prob{P} + \varepsilon^{100}\leq 3\tau^{1.5}\) with probability as high as \(\exp(-\poly(\varepsilon)\frac{d_1}{k})\) and hence this cannot occur for more than \(\exp(-\poly(\varepsilon)\frac{d_1}{k})\).

\medskip

    The second item follows directly from \pref{claim:edge-sampling}. The main effort is to show the first item. Assume that it doesn't hold. Define \(\mathcal{F}'_t\) by erasing \(f'_r\) for every \(r \in \bigcup_{g \in L}\supp_{\gamma+3\zeta+\varepsilon^{100}}(g)\) (and taking \(f'_r=f_r\) for the rest of the faces)\footnote{Formally, when we ``erase'' a function, we actually set the \(f'_r\) to be such that it agrees with no other functions. This can be done easily by extending the alphabet \(\Sigma\) to a large enough set.}. Then it still holds that 
     \[\agr_{\mathcal{D}}(\mathcal{F}'_t) \geq 2\tau^{1.5}.\]
     By the agreement soundness (assuming that \(\tau^{1.5}\geq \varepsilon_0\)), there is some \(h \in \overline{L}_{\tau^3,\eta}(t)\) such that \(A_{\eta}^{\mathcal{F}'}(h) \geq \tau^3\). On the other hand, there must by some \(g \in L\) such that \(\dist(h,g) \leq \gamma\). This implies that \(\prob{\supp_{\eta}^{\mathcal{F}'}(h) \setminus \supp_{\gamma+2\zeta}(g)} \leq \exp(-\poly(\zeta^2 k)) \ll \tau^3\). This is a contradiction since \(A_{\eta}^{\mathcal{F}'}(h) \subseteq \supp_{\eta}^{\mathcal{F}'}(h) \setminus \supp_{3\zeta+\gamma}(g)\) (as we erasing \(g\)'s support).
\end{proof}

\restateclaim{claim:good-tau-and-delta-for-every-t}

\begin{proof}[Proof of \pref{claim:good-tau-and-delta-for-every-t}]
    The proof logic is the following. We construct a sequence of lists \(L_1,L_2,\dots\) where \(L_1 \subseteq \bar{L}_{\tau_1,\zeta_1}(t)\), \(L_2 \subseteq \bar{L}_{\tau_2,\zeta_2}(t)\) and so on. If \(L_m \subseteq \bar{L}_{\tau_m,\zeta_m}(t)\) has the property required in the claim, i.e. that for \(1-\exp(-\poly(\varepsilon) \frac{d_1}{k})\) of the \(u \subseteq t\), \(L_m |_u\) is \(20\zeta_m\) dense inside \(\bar{L}_{\tau_{m+1},\zeta_m}(u)\) we conclude; else we continue to the next list \(L_{m+1}\). What we need to make sure is that there is an \(m \leq \frac{1}{\varepsilon^{10}}\) such that this process concludes with \(L_m\). To do so, we define a potential function \(p:\set{1,2,\dots, \frac{1}{\varepsilon^{10}}} \to [0,1]\) that is monotone non-decreasing (and \(p \leq 1\)), and show that if we don't conclude at step \(m\), then \(p(m+1) \geq p(m) + \varepsilon^{10}\), i.e.\ one of the first \(\frac{1}{\varepsilon^{10}}\) first lists must satisfy the property else \(p(\frac{1}{\varepsilon^{10}}) > 1\). The potential function is 
    \begin{equation} \label{eq:potential-function}
        p(m) = \Prob[r_1,r_2 \subseteq t, r_1,r_2 \sim D]{\bigcup_{g_j \in L_m} A_{9\zeta_{m-1}}(g_j)}.
    \end{equation}
    That is, \(p(m)\) measures how much agreement is explainable by \(L_m\), and we show that while the list \(L_m\) is not dense for most \(u \subseteq t\), then in fact we can find a list \(L_{m+1}\) whose functions `explains' more agreement. Details follow.
    
    Recall that \(\tau_m = \varepsilon^2 (1-\varepsilon^{10})^m\) and that \(\zeta_m = 3\eta \cdot 20^m\). Fix some \(t\). For every \(\tau_m,\zeta_m\) we choose some (arbitrary) \(L^{aux}_{m}(t) \subseteq \overline{L}_{\tau_{m}^3,\zeta_m}(t)\) that is \(3\zeta\)-dense and \(3\zeta\)-separated. We will sequentially construct lists \(L_m=\set{g_1,g_2,\dots,g_m} \subseteq \overline{L}_{\tau_{m},\zeta_{m}}(t)\), for \(m=1,2,\dots\) until \(L_m\) satisfies the claim. The potential function \(p\) shall be defined as in \pref{eq:potential-function} with respect to these lists, once we construct them. %

    For \(m=0\) we take \(L_0 =\emptyset \subseteq \overline{L}_{\tau_0,\zeta_0}(t)\). Let us describe how to construct \(L_m\) using \(L_{m-1}\), provided that \(L_{m-1}\) is not \(\zeta_m=20\zeta_{m-1}\)-dense in \(\overline{L}_{\tau_{m},\zeta_{m-1}}(u)\) for \(1-\exp(-\poly(\varepsilon)\frac{d_1}{k})\) of the \(u \in X(d_1)\):
        \begin{enumerate}
            \item Find some \(u \subseteq t\) such that:
            \begin{itemize}
                \item \(L_{m-1}\) is not \(20\zeta_{m-1}\)-dense in \(\overline{L}_{\tau_{m},\zeta_{m-1}}(u)\).
                \item \(u\) samples \(L^{aux}_{m} \cup L_{m-1}\) well.
                \item \(L^{aux}_{m}|_u\) is \(7\zeta_{m}\)-dense in \(\overline{L}_{\tau_{m},\zeta_{m-1}}(u)\). 
                \item For every \(g \in L^{aux}_{m}(t)\) it holds that \(\Abs{A_{9\zeta_{m-1}}^{\mathcal{F}_t}(g) - A_{9\zeta_{m-1}}^{\mathcal{F}_u}(g|_u) }\leq \varepsilon^{100}\).\footnote{The fraction of \(u \subseteq t\) such that these items do not hold is \(\exp(-\Omega(\poly(\varepsilon)\frac{d_1}{k}))\) by \pref{claim:edge-sampling} and \pref{claim:density-of-smaller-tau}. Hence if \(L_m\) is not \(20\zeta_{m-1}\)-dense for for \(1-\exp(-\Omega(\poly(\varepsilon)\frac{d_1}{k}))\) we can find a \(u \subseteq t\) such that this holds.}
            \end{itemize} 
            \item Find some \(h \in \overline{L}_{\tau_{m},\zeta_{m-1}}(u)\) that is \(20\zeta_{m-1}\)-far from every element in \(L_{m-1}\). There exists such an \(h\) since \(L_{m-1}\) is not \(20\zeta_{m-1}\)-dense in \(\overline{L}_{\tau_{m},\zeta_{m-1}}(u)\).
            \item Take some \(g=g_{m} \in L^{aux}_{m}\) such that \(\dist(g|_u,h_u) \leq 7\zeta_{m-1}\). Set \(L_{m} = L_{m-1} \cup \set{g_{m}}\). This step is possible since \(L^{aux}_{m}(t)|_u\) is \(7\zeta_{m}\)-dense inside \(\overline{L}_{\tau_{m},\zeta_{m}}(u)\).
        \end{enumerate}
    Observe that because \(\dist(g_{m}|_u,h) \leq 7\zeta_{m-1}\) then the fraction of \(r \in X(k)\) such that \(\dist(g_{m}|_r,h|_r) > 8\zeta_{m-1}\) has probability \(\exp(-\Omega(\zeta^2 k)) \ll \varepsilon^{100}\). In addition, when \(f_r \overset{1-\zeta_m}{\approx} h_r\) and \(h_r \overset{1-8\zeta_{m-1}}{\approx} g_{m}|_r\) then \(h_r \overset{1-9\zeta_{m-1}}{\approx} g_{m}\). Hence it holds that \(\prob{A_{9\zeta_{m-1}}(g_{m}|_u)} \geq \prob{A_{\zeta_{m-1}}(h)} - \varepsilon^{100}\). By assumption on \(u\) it holds that \(\prob{A_{9\zeta_{m-1}}(g_{m})} \geq \prob{A_{9\zeta_{m-1}}(g_{m}|_u)} -\varepsilon^{100}\) so
    \[\prob{A_{9\zeta_{m-1}}(g_{m})} \geq \prob{A_{9\zeta_{m-1}}(h)} - 2\varepsilon^{100} \geq \tau_{m-1}-2\varepsilon^{100}\geq \tau_{m}.\]
    That is, \(g_{m} \in \overline{L}_{\tau_{m},\zeta_{m}}(t)\) (and thus every \(L_m \subseteq \overline{L}_{\tau_{m},\zeta_{m}}(t)\)).

    Moreover, note that for every \(j \leq m\) it holds that \(\dist(g_{m}|_u,g_j|_u)\geq \dist(h,g_j|_u)-\dist(g_{m}|_u,h) \geq 13\zeta_{m}\). The face \(u\) samples \(L^{aux}_{m} \cup L_{m-1}\) well, hence \(\dist(g_{m},g_j) \geq \dist(g_{m}|_u,g_j|_u) - \zeta_{m-1} \geq 12\zeta_{m-1}\). This is much larger than \(9\zeta_{m-1} + \zeta_{j-1}\). Hence (as before) it holds that \(\prob{\supp_{9\zeta_{j-1}}(g_j) \cap \supp_{9\zeta_{m-1}}(g_{m})} \ll \varepsilon^{100}\). This implies that \(\prob{A_{9\zeta_{j-1}}(g_j) \setminus \bigcup_{j' \ne j} A_{9 \zeta_{j'-1}}(g_{j'})} \geq \frac{1}{2} \prob{A_{9\zeta_{j-1}}(g_j)}\).
    
    Thus it holds that while \(m \leq \varepsilon^{-30}\), then \(p(m) \geq \frac{1}{2}\sum_{j=1}^m \prob{A_{9\zeta_{j-1}}(g_j)} \geq \frac{m}{2}\varepsilon^{10}\). For \(m > 2\varepsilon^{10}\) this is greater than \(1\) hence the process must stop beforehand.
\end{proof}